\subsection{Semisimple endpoints and a small middle corner}

\begin{proposition}\label{C:semisimple-ends-small-middle}
Let $C$ be a finite category on objects $0,1,2$, with no arrows
from a larger object to a smaller one.  Suppose the two endpoint
monoid algebras are semisimple over $\mathbb Q$, and the middle
endomorphism monoid has at most three elements.  If
$H_*(BC;\mathbb Q)$ is bounded, then
\[
 H_4(BC;\mathbb Q)=0.
\]
In particular, such a category cannot have type $\mathbb CP^2$.
There is no bound on the endpoint orders or cross-hom cardinalities.
\end{proposition}
\begin{proof}
Put $R_i=\mathbb Q\operatorname{End}_C(i)$, and write
$S_{ij}=\mathbb Q\operatorname{Hom}_C(i,j)$, with its usual
left-$R_j$, right-$R_i$ actions.  Work throughout over
$\mathbb Q$.

The complete list of monoids of order at most three gives the
following elementary alternatives for $R_1$: it is semisimple;
it is the algebra of upper triangular two-by-two matrices or its
opposite; or it is
\[
 \mathbb Q\times D,\qquad D=\mathbb Q[\epsilon]/(\epsilon^2).
\]
For completeness, the two hereditary cases are an identity
adjoined to a two-element left-zero or right-zero band.  Their
difference of band elements spans the one-dimensional radical,
and the two primitive idempotents identify the algebra with the
displayed triangular algebra.  The nonhereditary case is
$\{1,a,z\}$ with $z$ a zero and $a^2=z$.  Its central
idempotent $z$ splits off the first factor, and $a-z$ is the
square-zero generator in the other factor.  The remaining
order-three cases are the cyclic group of order three, the
order-two group with a zero, that group with a new identity,
and the three-element semilattice chain; their rational
algebras are semisimple.  The order-one and order-two cases
are semisimple as well.

In every case both trivial $R_1$-modules have projective
dimension at most one.  In the dual-number case the trivial
module is precisely the projective factor supported at $z$.
Also the middle monoid has zero positive rational homology:
in the band or zero cases its classifying space is contractible,
and the other cases have finite-group homotopy type.

Filter the normalized nerve chains by the number of strict
changes of object.  That number is at most two.  The associated
graded terms are the ordinary corner bars for no change,
the two-sided correspondence bars for one change, and the
three-corner bar for the path $0\to1\to2$ for two changes.
The crossings contribute the horizontal degree $p$; internal
bar degree is $q$.  Exact endpoint coinvariants therefore give
\[
 E^1_{0,q}=0\quad(q>0),\qquad
 E^1_{1,q}=0\quad(q\geq2),\qquad
 E^1_{2,q}=\operatorname{Tor}^{R_1}_q(W,V),
\]
where
\[
 W=\mathbb Q\otimes_{R_2}S_{12},\qquad
 V=S_{01}\otimes_{R_0}\mathbb Q.
\]
The first equality uses the positive-homology assertion for
all three corners.  For the second equality, the direct
$0\to2$ correspondence has semisimple endpoints.  In each
other edge correspondence one endpoint is semisimple and the
other has trivial-module projective dimension at most one.
The formula remains valid if a cross hom-set is empty, in
which case the relevant module or bar term is zero.

For $q\geq2$ no spectral-sequence differential can enter
$(2,q)$.  Its first differential lands at $(1,q)$, which is
zero, and its possible second differential lands at
$(0,q+1)$, which is zero.  There are no larger horizontal
degrees.  There are also no other terms of total degree at
least four.  Consequently the filtration yields actual
vector-space isomorphisms
\begin{equation}\label{C:small-middle-high-homology}
 H_n(BC;\mathbb Q)
 \cong \operatorname{Tor}^{R_1}_{n-2}(W,V)
 \qquad(n\geq4).
\end{equation}
This does not assert that the lower-degree crossing
differentials vanish.

If $R_1$ is semisimple or hereditary, its Tor in degree two
and above is zero.  In the remaining case, only the
dual-number summands of $W,V$ can contribute positive Tor.
Every finite-dimensional $D$-module is a direct sum of free
modules and copies of $\mathbb Q=D/(\epsilon)$, by the
Jordan form of its square-zero operator.  If the respective
simple multiplicities are $b,a$, the elementary periodic
resolution with every differential multiplication by
$\epsilon$ gives
\[
 \operatorname{Tor}^D_j(W_D,V_D)\cong\mathbb Q^{ab}
 \qquad\text{for every }j\geq1.
\]
Thus a nonzero degree-two Tor would, by
\eqref{C:small-middle-high-homology}, give nonzero homology
in every degree at least four.  Bounded homology excludes
that possibility, proving the proposition.
\end{proof}

\begin{corollary}[A small middle corner with a zero]
\label{C:semisimple-ends-small-zero-middle}
The same conclusion holds if the middle monoid has a two-sided
zero and at most five elements, while the two endpoint algebras
remain semisimple over $\mathbb Q$.
\end{corollary}
\begin{proof}
Let $z$ be the middle zero.  The central idempotent $z$ gives
\[
 R_1=\mathbb Qz\times A,\qquad
 A=\mathbb QN/\mathbb Qz,
\]
where $A$ is the contracted algebra of the pointed finite
monoid $N$ and has dimension at most four.  The trivial
left and right $R_1$-modules are the projective factor
$\mathbb Qz$.  Thus the same strict-object filtration has
zero positive internal homology in horizontal degrees
zero and one.  It gives the stronger identification
\[
 H_n(BC;\mathbb Q)
 \cong\operatorname{Tor}^{R_1}_{n-2}(W,V)
 \cong\operatorname{Tor}^{A}_{n-2}(W_A,V_A)
 \qquad(n\geq3).
\]
The second isomorphism is in positive Tor degrees; the
semisimple $\mathbb Qz$ factor contributes none.
If degree-two Tor were nonzero, the contracted
four-dimensional persistence theorem
\ref{B:contracted-four-Tor} would make Tor nonzero in
arbitrarily high degrees.  The displayed identification
would contradict bounded rational homology of $BC$.
Consequently $H_4(BC;\mathbb Q)=0$.
\end{proof}

\begin{corollary}[Semisimple endpoints and two-point cross hom-sets]
\label{C:semisimple-ends-binary-homs}
Let $C$ have three objects $0,1,2$, no backward arrows,
semisimple rational algebras at $0$ and $2$, and
\[
 1\leq|C(0,1)|,\ |C(1,2)|,\ |C(0,2)|\leq2.
\]
The middle monoid is arbitrary.  Then $BC$ cannot have type
$\mathbb CP^2$.
\end{corollary}
\begin{proof}
Pass to the actual cross-act orbits at the two extreme
objects, replacing their corner monoids by trivial ones.
The two orbit-collapse functors induce rational homology,
and hence rational cohomology-ring, equivalences.  Indeed,
endpoint coinvariants are exact; the strict-object bar
filtration used above identifies this functor on each
associated-graded correspondence bar with that exact
coinvariant comparison.  Every resulting forward hom-set
is still nonempty and has at most two elements.

If an adjacent hom-set is a singleton, the singleton-cofinal
deletion removes the middle object; the result has two
trivial corners and is a graph.  Otherwise both adjacent
hom-sets have size two.  If the long hom-set has size one,
triangle compression gives a monoid with a $2\times2$
minimal band, whose rational degree-two products vanish
by Lemma~\ref{B:rank-one-square}.  If the long hom-set
also has size two, Corollary~\ref{C:middle-double-homs-rational}
gives nonzero first rational homology or zero second
rational homology.  Each alternative contradicts the
rational homology or cup square of a projective plane.
\end{proof}
